% References (paste into BibTeX / biblatex later if desired)
\begin{thebibliography}{99}
\sloppy

\bibitem{abadi2015}
Abadi, M., A. Agarwal, P. Barham, E. Brevdo, Z. Chen, C. Citro, G. S. Corrado, A. Davis, J. Dean, M. Devin, S. Ghemawat, I. Goodfellow, A. Harp, G. Irving, M. Isard, Y. Jia, R. Jozefowicz, L. Kaiser, M. Kudlur, J. Levenberg, D. Mane, R. Monga, S. Moore, D. Murray, C. Olah, M. Schuster, J. Shlens, B. Steiner, I. Sutskever, K. Talwar, P. Tucker, V. Vanhoucke, V. Vasudevan, F. Viegas, O. Vinyals, P. Warden, M. Wattenberg, M. Wicke, Y. Yu, and X. Zheng (2015). TensorFlow: Large-scale machine learning on heterogeneous systems. Software available from tensorflow.org.

\bibitem{ali2023}
Ali, M. (2023). Distributed processing using ray framework in python.

\bibitem{amdahl1967}
Amdahl, G. M. (1967). Validity of the single processor approach to achieving large scale computing capabilities. In \textit{Proceedings of the April 18--20, 1967, Spring Joint Computer Conference, AFIPS '67 (Spring)}, New York, NY, USA, pp. 483--485. Association for Computing Machinery.

\bibitem{bates2008}
Bates, D. and D. Watts (2008). \textit{Nonlinear Regression Analysis and Its Applications}, pp. 32--66.

\bibitem{chen2024ieee}
Chen, Y., O. M. Saad, A. Savvaidis, F. Zhang, Y. Chen, D. Huang, H. Li, and F. Aziz Zanjani (2024). Deep learning for p-wave first-motion polarity determination and its application in focal mechanism inversion. \textit{IEEE Transactions on Geoscience and Remote Sensing} \textbf{62}, 1--11.

\bibitem{chen2024ess}
Chen, Y., Savvaidis, A., Siervo, D., Huang, D., \& Saad, O. M. (2024). Near real-time earthquake monitoring in Texas using the highly precise deep learning phase picker. \textit{Earth and Space Science}, \textbf{11}, e2024EA003890. \url{https://doi.org/10.1029/2024EA003890}

\bibitem{choi2025}
Choi, S. B. (2025). Design and experimental research of on-device style transfer models based on PyTorch to ONNX. \textit{Scientific Reports}, 15, Article number: 9345.

\bibitem{czyzewski2021}
Czyzewski, M. A. (2021). Transfer Learning Between Different Architectures Via Weights Injection. arXiv preprint arXiv:2101.02757.

\bibitem{feng2022}
Feng, T., S. Mohanna, and L. Meng (2022). Edgephase: A deep learning model for multi-station seismic phase picking. \textit{Geochemistry, Geophysics, Geosystems} \textbf{23}(11), e2022GC010453.

\bibitem{friedman2001}
Friedman, J. H. (2001). Greedy function approximation: a gradient boosting machine. \textit{The Annals of Statistics} \textbf{29}(5), 1189--1232.

\bibitem{gustafson1988}
Gustafson, J. L. (1988). Reevaluating amdahl's law. \textit{Commun. ACM} \textbf{31}(5), 532--533.

\bibitem{herlihy2012}
Herlihy, M. and N. Shavit (2012). \textit{The Art of Multiprocessor Programming, Revised Reprint}. Morgan Kaufmann.

\bibitem{johnston1997}
Johnston, J. and J. DiNardo (1997). \textit{Econometric Methods} (4th ed.). McGraw-Hill.

\bibitem{lim2025}
Lim, C. S. Y., S. Lapins, M. Segou, and M. J. Werner (2025). Deep learning phase pickers: how well can existing models detect hydraulic-fracturing induced microseismicity from a borehole array? \textit{Geophysical Journal International} \textbf{240}(1), 535--549. \url{https://doi.org/10.1093/gji/ggae386}

\bibitem{mccool2012}
McCool, M., J. Reinders, and A. Robison (2012). \textit{Structured Parallel Programming: Patterns for Efficient Computation}. Elsevier Science.

\bibitem{moritz2018}
Moritz, P., R. Nishihara, S. Wang, A. Tumanov, R. Liaw, X. Liang, M. Elibol, Z. Yang, W. Paul, M. I. Jordan, and I. Stoica (2018). Ray: A Distributed Framework for Emerging AI Applications. In \textit{13th USENIX Symposium on Operating Systems Design and Implementation (OSDI 18)}, Carlsbad, CA, pp. 561--577.

\bibitem{mousavi2020}
Mousavi, S., W. Ellsworth, Z. Weiqiang, L. Chuang, and G. Beroza (2020). Earthquake transformer---an attentive deep-learning model for simultaneous earthquake detection and phase picking. \textit{Nature Communications} \textbf{11}, 3952.

\bibitem{mousavi2022}
Mousavi, S. M. and G. C. Beroza (2022). Deep-learning seismology. \textit{Science} \textbf{377}(6607), eabm4470.

\bibitem{munoz2025}
Muñoz, C., V. Salles, C. Skevofilax, and A. Savvaidis (2025). SCMLPick: A SeisComP Module Implementing Machine Learning Phase Picking for Real-Time Seismic Monitoring. [Manuscript submitted for publication]. Texas Seismological Network, University of Texas at Austin.

\bibitem{muenchmeyer2022}
Münchmeyer, J., J. Woollam, A. Rietbrock, F. Tilmann, D. Lange, T. Bornstein, T. Diehl, C. Giunchi, F. Haslinger, D. Jozinović, A. Michelini, J. Saul, and H. Soto (2022). Which picker fits my data? a quantitative evaluation of deep learning based seismic pickers. \textit{Journal of Geophysical Research: Solid Earth} \textbf{127}.

\bibitem{onnx2021}
ONNX Community (2021). ONNX: Open Neural Network Exchange. Available at \url{https://onnx.ai/}

\bibitem{parida2025}
Parida, S. K., and other authors (2025). How Do Model Export Formats Impact the Development of ML-Enabled Systems? A Case Study on Model Integration. arXiv preprint arXiv:2502.00429.

\bibitem{paszke2019}
Paszke, A., S. Gross, F. Massa, A. Lerer, J. Bradbury, G. Chanan, T. Killeen, Z. Lin, N. Gimelshein, L. Antiga, A. Desmaison, A. Köpf, E. Yang, Z. DeVito, M. Raison, A. Tejani, S. Chilamkurthy, B. Steiner, L. Fang, J. Bai, and S. Chintala (2019). PyTorch: An imperative style, high-performance deep learning library. \textit{Advances in Neural Information Processing Systems} 32, 8024--8035.

\bibitem{raydocs}
Ray.io (n.d.). What is ray core? \url{https://docs.ray.io/en/latest/ray-core/walkthrough.html}

\bibitem{saad2022caps}
Saad, O. M. and Y. Chen (2022). Capsphase: Capsule neural network for seismic phase classification and picking. \textit{IEEE Transactions on Geoscience and Remote Sensing} \textbf{60}, 1--11.

\bibitem{saad2022unsupervised}
Saad, O. M., Y. Chen, A. Savvaidis, W. Chen, F. Zhang, and Y. Chen (2022). Unsupervised deep learning for single-channel earthquake data denoising and its applications in event detection and fully automatic location. \textit{IEEE Transactions on Geoscience and Remote Sensing} \textbf{60}, 1--10.

\bibitem{saad2022jgr}
Saad, O. M., Y. Chen, A. Savvaidis, S. Fomel, and Y. Chen (2022). Real-time earthquake detection and magnitude estimation using vision transformer. \textit{Journal of Geophysical Research: Solid Earth} \textbf{127}(5), e2021JB023657.

\bibitem{saad2023}
Saad, O. M., Y. Chen, D. Siervo, F. Zhang, A. Savvaidis, G.-c. D. Huang, N. Igonin, S. Fomel, and Y. Chen (2023). Eqcct: A production-ready earthquake detection and phase-picking method using the compact convolutional transformer. \textit{IEEE Transactions on Geoscience and Remote Sensing} \textbf{61}, 1--15.

\bibitem{saad2021}
Saad, O. M., G. Huang, Y. Chen, A. Savvaidis, S. Fomel, N. Pham, and Y. Chen (2021). Scalodeep: A highly generalized deep learning framework for real-time earthquake detection. \textit{Journal of Geophysical Research: Solid Earth} \textbf{126}(4), e2020JB021473.

\bibitem{savvaidis2019}
Savvaidis, A., B. Young, G. D. Huang, and A. Lomax (2019). Texnet: A statewide seismological network in texas. \textit{Seismological Research Letters} \textbf{90}(4), 1702--1715.

\bibitem{secondtalent}
SecondTalent (n.d.). PyTorch vs. TensorFlow: usage, popularity, and performance. \url{https://www.secondtalent.com/resources/pytorch-vs-tensorflow-usage-popularity-and-performance/}

\bibitem{sheen2023}
Sheen, D.-H., Y.-J. Seong, S. Jung, and J.-h. Park (2023). Real-time application of deep learning seismic phase picker to earthquake monitoring. Poster S31G-0424, \textit{AGU Fall Meeting}, San Francisco, CA.

\bibitem{skevofilax2026}
Skevofilax, C., and A. Savvaidis (2026). RAPID: A Generalized High-Performance Parallelization Framework for Real-Time Deep Learning Seismic Phase Picking. [Unpublished manuscript]. Texas Seismological Network, University of Texas at Austin.

\bibitem{tan2021}
Tan, Y. J., F. Waldhauser, W. Ellsworth, M. Zhang, Z. Weiqiang, M. Michele, L. Chiaraluce, G. Beroza, and M. Segou (2021). Machine-learning-based high-resolution earthquake catalog reveals how complex fault structures were activated during the 2016--2017 central italy sequence. \textit{The Seismic Record} \textbf{1}, 11--19.

\bibitem{verma2024}
Verma, G., S. Raskar, M. Emani, and B. Chapman (2024). Cross-Feature Transfer Learning for Efficient Tensor Program Generation. \textit{Applied Sciences}, 14(2), 513.

\bibitem{wang2022}
Wang, H., J. Gu, and Y. Li (2022). An Empirical Study of Challenges in Converting Deep Learning Models. In \textit{Proceedings of the 31st ACM SIGSOFT International Symposium on Software Testing and Analysis (ISSTA '22)}, 206--217.

\bibitem{wang2021}
Wang, T., D. Trugman, and Y. Lin (2021). Seismogen: Seismic waveform synthesis using gan with application to seismic data augmentation. \textit{Journal of Geophysical Research: Solid Earth} \textbf{126}(4), e2020JB020077.

\bibitem{woollam2022}
Woollam, J., J. Münchmeyer, F. Tilmann, A. Rietbrock, D. Lange, T. Bornstein, T. Diehl, C. Giunchi, F. Haslinger, and D. Jozinović (2022). SeisBench---A Toolbox for Machine Learning in Seismology. \textit{Seismological Research Letters} \textbf{93}(3), 1695--1709.

\bibitem{xiao2021}
Xiao, Z., J. Wang, C. Liu, J. Li, L. Zhao, and Z. Yao (2021). Siamese earthquake transformer: A pair-input deep-learning model for earthquake detection and phase picking on a seismic arrival-time picking method. \textit{Journal of Geophysical Research: Solid Earth} \textbf{126}(5), e2020JB021444.

\bibitem{yeck2020}
Yeck, W. L., J. M. Patton, Z. E. Ross, G. P. Hayes, M. R. Guy, N. B. Ambruz, D. R. Shelly, H. M. Benz, and P. S. Earle (2020). Leveraging Deep Learning in Global 24/7 Real-Time Earthquake Monitoring at the National Earthquake Information Center, \textit{Seismol. Res. Lett.} \textbf{92}, 469--480, doi: 10.1785/0220200178.

\bibitem{yu2024}
Yu, Z., Y. Jiang, X. Jing, and H. Zheng (2024). Study on geomagnetic observations associated with three major earthquakes in southwest china through a novel deep learning framework. \textit{IEEE Transactions on Geoscience and Remote Sensing} \textbf{62}, 1--18.

\bibitem{zhu2018}
Zhu, W. and G. C. Beroza (2018). Phasenet: a deep-neural-network-based seismic arrival-time picking method. \textit{Geophysical Journal International} \textbf{216}(1), 261--273.

\end{thebibliography}
